% Isolated source-corrective candidate; preserve all existing statement labels.
\begin{proposition}
\label{prop:ch20-curve-laplacian}
Let \(n=1\), let \(\mathcal X\) be a regular model of the smooth projective
curve \(X\), and write \(\mathcal X_s=\sum_{i\in I}b_iE_i\).  Let
\(D=\sum_i a_iE_i\) be a vertical real divisor, with associated
vertex values \(u_i=\varphi_D(x_{E_i})=a_i/b_i\).  Then
\begin{equation}
\label{eq:ch20-curve-laplacian}
 (D\cdot E_i)=\sum_{j\ne i}b_j(E_j\cdot E_i)(u_j-u_i),
 \qquad \sum_i b_i(D\cdot E_i)=0.
\end{equation}
The off-diagonal coefficients \(b_j(E_j\cdot E_i)\) are nonnegative.
Multiplying the first identity by \(b_i\) gives symmetric edge
coefficients \(b_i b_j(E_i\cdot E_j)\).

Suppose moreover that \(\mathcal X\) is regular semistable, with reduced
special fibre and skeleton \(\Gamma=\Delta_{\mathcal X}\).
Write \(k=\widetilde K\).  To the edge corresponding to a closed node
\(z_e\) assign the residue weight \(c_e=[\kappa(z_e):k]\).
Its modulus length is one under \(-\log|\pi|=1\); the local equation
\(xy=\pi\) is understood on an etale node chart, and does not assert
that \(z_e\) is \(k\)-rational.  With the weighted outgoing-slope
operator
\[
 \Delta^{(c)}_\Gamma F
   :=\sum_{x\in\Gamma}
       \left(\sum_{\nu\in T_x\Gamma}
                   c_{e(\nu)}\partial_\nu F(x)\right)\delta_x,
\]
one has
\begin{equation}
\label{eq:ch20-curve-laplacian-graph}
 \sum_i(D\cdot E_i)\delta_{x_{E_i}}
                 =\Delta^{(c)}_\Gamma\varphi_D.
\end{equation}
The sign is the outgoing-slope sign of \eqref{eq:graph-laplacian}:
a convex corner has positive mass.  The operator itself agrees with
that unweighted formula only when every node has residue degree one.
\end{proposition}

\begin{proof}
Write \(D=\sum_j b_j u_jE_j\).  The special fibre is the principal
divisor of \(\pi\), so \(\sum_j b_j(E_j\cdot E_i)=0\).
Subtracting \(u_i\) times this equality from
\(D\cdot E_i=\sum_j b_j u_j(E_j\cdot E_i)\) proves the first
identity in \eqref{eq:ch20-curve-laplacian}; its \(j=i\) term is zero.
The second is \(D\cdot\mathcal X_s=0\), or follows by multiplying
by \(b_i\), summing, and cancelling terms using symmetry of the
intersection pairing.  For distinct prime divisors the local
intersection multiplicities are nonnegative, which proves the
off-diagonal assertion.  All these equalities extend real-linearly.

In the semistable case the intersection at a node is transverse.
Its local intersection length is one over \(\kappa(z_e)\); its
degree over \(k\) is therefore \([\kappa(z_e):k]\).  Consequently
\begin{equation}
\label{eq:ch20-node-intersection-residue-degree}
             (E_i\cdot E_j)
        =\sum_{z\in E_i\cap E_j}[\kappa(z):k]
                   \qquad(i\ne j).
\end{equation}
The right side is a degree-weighted sum, not the number of closed nodes.
The length of each node annulus is
\(-\log|\pi|=1\), computed on its etale chart.  The affine model
function has endpoint values \(u_i,u_j\), so its outgoing slope
at \(x_{E_i}\) on that edge is \(u_j-u_i\).  Summing with the
node weights gives
\[
 \sum_{e\text{ incident to }i}c_e\partial_e\varphi_D(x_{E_i})
    =\sum_{j\ne i}(E_i\cdot E_j)(u_j-u_i)
    =D\cdot E_i.
\]
The first equality is \eqref{eq:ch20-node-intersection-residue-degree},
and the second is \eqref{eq:ch20-curve-laplacian} with \(b_i=1\).
There are no interior slope breaks, so this is the complete measure
identity \eqref{eq:ch20-curve-laplacian-graph}.
The residue-degree convention is also that of
\cite[Remark~1.2.4 and \S2.2.3]{Thu05}; it agrees with the later
arbitrary-field formula \eqref{eq:ch31-graph-laplacian}.
\end{proof}

